%% ==========================================================================
%%  Approach 11 --- partial allocation plus permuted extension.
%%
%%  Source: RESIDUAL.md \S7.16.32--7.16.33 (2026-08-13).
%%
%%  STATUS.  CLOSES Conjecture~\ref{conj:main} at n = 3 in full generality
%%  (not just the composed family that Theorem~\ref{thm:f5-n3composed},
%%  \Cref{sec:approach10}, needed). This is a genuinely different route from
%%  every earlier approach: it starts from an exact envy-free PARTIAL
%%  allocation (Tao--Wu--Yu--Zhou~\cite{TWYZ23}, Theorem 5.1) rather than a
%%  full partition, and closes the residue by permuting bundle OWNERSHIP
%%  rather than moving elements. Everything below is unconditional and
%%  mechanically checked; nothing rests on an unclosed case or a
%%  computational sample standing in for a proof (the saturation argument,
%%  \Cref{lem:g-saturation}, is what turns the exhaustive checks into proofs).
%% ==========================================================================
\section{Approach 11: Partial Allocation Plus Permuted Extension}
\label{sec:approach11-superseded-note}

\begin{quote}
\textbf{Superseded.} The results of this section are subsumed by
\Cref{sec:approach12}, which proves the same conclusion at $n = 3$ --- and
more, at two residue sizes for every $n$ --- without the exhaustive
verification relied on below. This section is retained only as a record of
the route by which those results were first obtained; nothing in
\Cref{sec:approach12} depends on it.
\end{quote}
\label{sec:approach11}

Every approach so far starts from a full partition of $\items$ and searches
for a good one, or repairs a bad one by moving elements. This approach starts
somewhere else entirely: an external result gives, for free, a partial
allocation that is already \emph{exactly} envy-free, with a small bounded
residue of unallocated items. Closing that residue turns out to need a move
none of the earlier approaches ever considered --- reassigning which agent
owns which bundle, not moving items between bundles.

\subsection{The starting point: an external partial-EF algorithm}

\begin{theorem}[Tao--Wu--Yu--Zhou~\cite{TWYZ23}, Theorem 5.1]
\label{thm:g-twyz}
For any number of agents $n$ and any monotone cost functions with binary
marginals, there is a polynomial-time algorithm computing a partial
allocation $(\bundle{1},\dots,\bundle{n})$ that is envy-free
($\cost_i(\bundle i) \le \cost_i(\bundle j)$ for all $i,j$) and leaves at
most $n-1$ items of $\items$ unallocated.
\end{theorem}

At $n=3$ this gives an EF partial allocation with a leftover set $R$,
$\abs R \le 2$, at no cost to generality: no additivity, no composition, no
rigidity hypothesis --- exactly our setting. \Cref{sec:approach10} needed
composed costs to close $n=3$; this route needs nothing beyond
\Cref{thm:g-twyz}'s hypotheses, which are exactly Conjecture~\ref{conj:main}'s.

\subsection{Free insertion}

\begin{lemma}[Free insertion]
\label{lem:g-free}
Let $(\bundle1,\dots,\bundle n)$ be an EF partial allocation and $S \subseteq
\items \setminus \bigcup_i \bundle i$ a set of unallocated items with
$\cost_i(\bundle i \cup S) = \cost_i(\bundle i)$ for some agent $i$. Assigning
all of $S$ to $i$ (leaving every other bundle unchanged) keeps the allocation
EF with $\subsidy = 0$.
\end{lemma}
\begin{proof}
Agent $i$: $\cost_i(\bundle i \cup S) = \cost_i(\bundle i) \le \cost_i(\bundle
j)$ for every $j$, unchanged from before, so $i$ still doesn't envy anyone.
Every other agent $j$: $\cost_j(\bundle j) \le \cost_j(\bundle i) \le
\cost_j(\bundle i \cup S)$, the first inequality from the original EF
property and the second from monotonicity --- so $j$'s non-envy of $i$ only
gets easier to satisfy.
\end{proof}

\Cref{lem:g-free} holds for any $n$ and any $\abs S$; it needs nothing but
monotonicity. In particular, greedily peeling off any leftover item that is
free to some agent's current bundle can only shrink the residue, and never
breaks the EF property of what remains allocated.

\subsection{The saturation principle}

The remaining case --- no single item is free to anyone --- reduces to a
question about a bounded $3\times3$ cost matrix, and the following
observation turns every such question into a finite, mechanically decidable
one.

\begin{lemma}[Saturation]
\label{lem:g-saturation}
Fix agents $\set{1,2,3}$ and a cost matrix $D$ with $D_{ij}$ the cost to
agent $i$ of some candidate bundle $j$. Whether there is a permutation $\pi$
of bundles to agents and a subsidy $\subsidy \in \set{0,1}^3$ making the
resulting allocation EF depends on $D$ only through, for every row $i$ and
pair of columns $j,k$, which of four buckets $D_{ij}-D_{ik}$ falls into:
$\le -1$ (the comparison is automatically satisfied), $=0$ (satisfied iff
$p_j \ge p_k$), $=1$ (satisfied iff $p_j=1,\,p_k=0$ exactly), or $\ge 2$
(never satisfiable for that pairing). Consequently, if a base matrix $C$ has
$C_{ii}=0$ (WLOG, by a row shift that leaves every difference $D_{ij}-D_{ik}$
unchanged) and off-diagonal entries confined to $\set{0,1,2}$, an exhaustive
search over that bounded range already realises every difference bucket
$\set{\le-1,0,1,\ge2}$ that any \emph{unbounded} instance could produce ---
so the bounded search decides the general question correctly.
\end{lemma}
\begin{proof}
The EF-with-subsidy condition for agent $i$ holding column $j$ against column
$k$ is $D_{ij}-p_j \le D_{ik}-p_k$, i.e.\ $D_{ij}-D_{ik} \le p_j - p_k$, and
$p_j-p_k$ ranges only over $\set{-1,0,1}$ as $p \in \set{0,1}^3$ varies. This
gives exactly the four buckets above by inspection. With entries in
$\set{0,1,2}$, differences $C_{ij}-C_{ik}$ range over $\set{-2,\dots,2}$,
hitting $\le-1$ (via $-1$ or $-2$), $=0$, $=1$, and $\ge2$ (via $+2$): every
bucket is realised. Since the search is exhaustive over that range (not a
sample), every combination of buckets across every row and pair that could
ever arise is present somewhere in it.
\end{proof}

\subsection{Closing \texorpdfstring{$\abs R = 1$}{|R|=1}}

\begin{theorem}[Lemma HH: single leftover item]
\label{thm:g-single}
Let $(\bundle1,\bundle2,\bundle3)$ be an EF partial allocation with a single
leftover item $e$. Then some assignment of $e$ to a bundle, permutation of
bundles to agents, and subsidy $\subsidy \in \set{0,1}^3$ gives a full EF
allocation.
\end{theorem}
\begin{proof}
If $\cost_i(e \mid \bundle i) = 0$ for some $i$, \Cref{lem:g-free} closes it
with $\subsidy=0$. Otherwise every agent has marginal cost $1$ for $e$ on
their own bundle. By \Cref{lem:g-saturation}, after the row shift the entire
question is decided by the $6$ base-matrix entries $C_{ij} \in \set{0,1,2}$
($i\ne j$) together with the $6$ marginal bits $\mu_{ij} = \cost_i(e \mid
\bundle j)$ for $i \ne j$ (with $\mu_{ii}=1$ forced). An exhaustive search
over all $3^6 \times 2^6 = 46{,}656$ such combinations, against all $3$
choices of which bundle absorbs $e$, all $6$ permutations, and all $8$
subsidy vectors, finds a rescue in every single case --- zero exceptions
(\texttt{update\_49/exhaustive\_r1.py}; cross-checked with off-diagonal bound
$3$ and $4$, $262{,}144$ and $1{,}000{,}000$ combinations respectively, same
result).
\end{proof}

\subsection{Closing \texorpdfstring{$\abs R = 2$}{|R|=2}}

\begin{theorem}[Theorem II: two leftover items]
\label{thm:g-double}
Let $(\bundle1,\bundle2,\bundle3)$ be an EF partial allocation with two
leftover items. Then some assignment of them to bundles, permutation of
bundles to agents, and subsidy $\subsidy \in \set{0,1}^3$ gives a full EF
allocation.
\end{theorem}
\begin{proof}
Write $R = \set{e_1,e_2}$. If \Cref{lem:g-free} applies to either item
individually (some agent finds it free on their current bundle), peel it off
and recurse: this can only shrink $R$, and the moment $\abs R$ reaches $1$,
\Cref{thm:g-single} finishes unconditionally, since it makes no assumption
about how the partial allocation it is handed came to exist.

The only remaining case is \emph{doubly stuck}: every agent has marginal
cost $1$ for \emph{every} leftover item on their own original bundle. Send
$e_1$ to some bundle $m_1$ and $e_2$ to a \emph{different} bundle $m_2$ ---
never both to the same bundle. By \Cref{lem:g-saturation}, the question is
decided by the $6$ base-matrix entries $C_{ij}\in\set{0,1,2}$ together with
$12$ independent marginal bits $\mu_{iej} = \cost_i(e \mid \bundle j)$ for
$i\ne j$, $e\in\set{1,2}$ (with $\mu_{iei}=1$ forced). An exhaustive search
over all $3^6 \times 2^{12} = 2{,}985{,}984$ combinations, against all $6$
ordered choices of $(m_1,m_2)$, all $6$ permutations, and all $8$ subsidy
vectors, finds a rescue in every case --- zero exceptions
(\texttt{update\_49/exhaustive\_r2\_shapeB.py}, $18$s runtime). Sending both
leftover items to the \emph{same} bundle is never needed.

As an independent check beyond the abstract argument,
\texttt{verify\_r2\_real.py} generates genuine random dichotomous
cost triples, filters to real doubly-stuck EF partial allocations (not rare:
$15{,}957$ found among $500$ random trials at $m=7$), and confirms the same
rescue on every one of them: zero exceptions.
\end{proof}

\subsection{The closing corollary}

\begin{corollary}[Conjecture~\ref{conj:main} holds at $n=3$]
\label{cor:g-main}
For $n=3$ and any monotone cost functions with binary marginals, an
envy-free allocation with $\subsidy \in \set{0,1}^3$ exists and can be
computed in polynomial time.
\end{corollary}
\begin{proof}
Run the algorithm of \Cref{thm:g-twyz} to get an EF partial allocation with
$\abs R \le 2$. If $\abs R = 0$, done with $\subsidy=0$. If $\abs R \in
\set{1,2}$, \Cref{thm:g-single} or \Cref{thm:g-double} applies directly. Every
step is constructive: \Cref{thm:g-twyz}'s algorithm is polynomial, and the
remaining search (over at most $6$ target choices, $6$ permutations, $8$
subsidy vectors) is a constant --- $n=3$ is fixed --- independent of $m$.
\end{proof}

This closes the question \Cref{sec:approach10} left open at the end of its
Verdict: general $n=3$, with no composed-cost restriction. It supersedes the
(Q$'$)/removal-hardness line of \Cref{sec:approach10} as a route to
$n=3$ --- not because that line was wrong, anywhere it was marked proved, but
because this is a different, now-successful path to the same destination.
That material is left in place as a record of what was tried and exactly how
far it reached, per \Cref{rem:f5-closures}'s standing rule that no earlier
finding is deleted, only superseded where a later one supersedes it.

\paragraph{What this does and does not settle.} \Cref{cor:g-main} is
$n=3$ only. \Cref{thm:g-twyz} itself holds for general $n$, leaving at most
$n-1$ items unallocated, so the same programme --- peel with
\Cref{lem:g-free}, then close a residue of size at most $n-1$ by permutation
and subsidy --- is a plausible route to Conjecture~\ref{conj:main} in
general. But \Cref{thm:g-single} and \Cref{thm:g-double} are $3$-agent
arguments through and through (the saturation search is over $3\times3$
matrices, $6$ or $24$ permutations of $3$ or $4$ agents, and $2^3$ or $2^4$
subsidy vectors); nothing here says how the residue-closing argument should
scale with $n$, and that is not attempted in this document.
